\section{Introduction}
\label{sec:introduction}

As Kubernetes has consolidated as the \textit{de facto} operating system for cloud-native microservices~\cite{gannon_2017, kosinska_2023}, organizations increasingly rely on Chaos Engineering to proactively validate fault tolerance, self-healing controllers, and observability pipelines prior to production outages~\cite{basiri_2016, owotogbe2026tenyears}. Moving beyond first-generation virtual machine terminators such as Netflix's Chaos Monkey~\cite{basiri_2016}, the Cloud Native Computing Foundation (CNCF) has fostered two incubating Kubernetes-native chaos engineering frameworks: \textbf{Chaos Mesh} (originally developed at PingCAP) and \textbf{LitmusChaos} (originally developed at MayaData/Harness)~\cite{owotogbe2026tenyears, onkalo2025improving}. Large-scale repository mining by Owotogbe et al.~\cite{owotogbe2026chaos} across $1{,}275$ open-source projects reveals that Chaos Mesh ($24.24\%$ of repositories) and LitmusChaos ($10.35\%$) are by far the two most widely adopted Kubernetes-native fault injection platforms in industry and academia.

Despite both tools exposing declarative Kubernetes Custom Resource Definitions (CRDs) to inject network, pod, and resource stress faults, their underlying architectural philosophies differ fundamentally. Chaos Mesh employs a \textit{DaemonSet-centric direct injection model}: a centralized controller dispatches gRPC commands directly to a node-resident privileged \texttt{chaos-daemon} pod, which enters the target container's Linux namespaces (\texttt{nsenter}) to configure \texttt{tc-netem} traffic queues, \texttt{iptables} rules, or \texttt{stress-ng} workers in place. In contrast, LitmusChaos employs an \textit{ephemeral runner orchestration model}: for each \texttt{ChaosEngine} workflow, the \texttt{chaos-operator} dynamically spawns a \texttt{chaos-runner} pod, which in turn schedules transient Kubernetes Job pods (\texttt{go-runner}) on the target node to execute the fault and evaluate declarative resilience probes (\texttt{httpProbe}, \texttt{cmdProbe}, \texttt{k8sProbe}, \texttt{promProbe}), storing results and Resilience Scores in a persistent control plane (\texttt{ChaosCenter} and MongoDB)~\cite{onkalo2025improving, opara2025chaos}.

While recent reviews by Onkalo~\cite{onkalo2025improving}, Mailewa et al.~\cite{mailewa2025review}, Opara et al.~\cite{opara2025chaos}, and Ambani~\cite{ambani2025systematic} compare the feature matrices and documentation of Chaos Mesh and LitmusChaos qualitatively, they explicitly highlight the absence of controlled empirical benchmarks evaluating their runtime behavior head-to-head on an identical Kubernetes testbed. Specifically, engineering teams lack empirical evidence regarding how ephemeral runner pod creation versus resident DaemonSet execution affects: (i) fault actuation delay and temporal precision, (ii) cluster CPU/memory overhead in resource-constrained environments, (iii) telemetry interference and diagnostic reproducibility in multimodal observability stacks, and (iv) autonomous steady-state hypothesis verification.

To bridge this gap, we conduct a controlled experiment following the guidelines of Wohlin et al.~\cite{wohlin_2012} and the Goal/Question/Metric (GQM) methodology~\cite{basili_1994}, evaluating Chaos Mesh and LitmusChaos across $N=30$ paired, stratified fault campaigns on the CNCF OpenTelemetry Astronomy Shop benchmark~\cite{opentelemetry_demo, wang2024building} monitored by a correlated multimodal observability stack (OpenTelemetry, Prometheus, Loki, and Tempo). We investigate four Research Questions (RQs):
\begin{itemize}[leftmargin=1.2em, itemsep=2pt, topsep=3pt]
    \item \textbf{RQ01 (Temporal Actuation and Teardown Precision):} How do the execution architectures of Chaos Mesh (resident DaemonSet) and LitmusChaos (ephemeral runner pods) impact fault actuation latency ($\Delta T_{\text{actuate}}$) and post-experiment cleanup latency ($\Delta T_{\text{teardown}}$)?
    \item \textbf{RQ02 (Control-Plane and Active Resource Overhead):} What is the steady-state and active-campaign CPU ($\Delta_{\text{CPU}}$) and memory ($\Delta_{\text{RAM}}$) footprint imposed by Chaos Mesh versus LitmusChaos?
    \item \textbf{RQ03 (Diagnostic Reproducibility and Telemetry Noise):} Do equivalent fault specifications in Chaos Mesh and LitmusChaos yield consistent root-cause diagnostic accuracy ($Acc$) and convergence latency ($MTTDiag$) under automated multimodal Root-Cause Analysis (RCA), and how much auxiliary telemetry noise is introduced by transient runner pods?
    \item \textbf{RQ04 (Native Steady-State Verification Expressiveness):} How effectively do LitmusChaos's native multi-type probes (\texttt{promProbe}, \texttt{k8sProbe}, \texttt{httpProbe}, \texttt{cmdProbe}) detect process-visible and process-invisible SLO violations compared to Chaos Mesh's native \texttt{StatusCheck}?
\end{itemize}
